Across the 226/300 gated questions, the intervention does \emph{not} inject
hedging vocabulary --- hedge-marker frequency \emph{drops} ($0.57\to0.27$ per
trace) --- and traces get shorter ($95\to78$ words). What changes is the mode
of reasoning: assertive shortcut arguments are replaced by definitional,
evidence-first ones, and the final answer changes on 83/226 gated questions
(13 overconfident-wrong become \emph{correct} under M4; 32/48 gated
confident-right survive intact). Ablating the confidence projection during
the \texttt{think} phase appears to remove a ``commit early'' pressure rather
than adding doubt. Three representative pairs (traces truncated):

\noindent\textbf{Overconfident-wrong $\to$ corrected} (\texttt{mmlu-4592}; gold C
, M4 state overconfident wrong $\to$ nonconfident wrong). The baseline trace jumps from a vague association to an answer; the steered trace states the definition and derives the (gold) answer.
\begin{quote}\scriptsize\emph{Question:} If the price of firm A's cell phone service rises by 5 percent and the quantity demanded for firm B's cell phone service increases by 10 percent, we can say that Options: A. demand for firm B is price elastic. B. supply\\[2pt]
\textbf{before:} \texttt{\textless{}think\textgreater{}The question asks about the relationship between the price of firm A's service and the quantity demanded for firm B's service. We know that the price of firm A's service increased and the quantity demanded for firm B's service increased. This suggests that the demand for firm B's service is likely to be price elastic. \textless{}/think\textgreater{} \textbackslash\{\}boxed\{A\}} $\to$ \texttt{A}\\[2pt]
\textbf{after:} \texttt{\textless{}think\textgreater{}The cross-price elasticity of demand measures the percentage change in the quantity demanded of one good in response to a percentage change in the price of another good. Since the price of firm A's service increased and the quantity demanded for firm B increased, the cross-price elasticity is positive. We can conclude that the two goods are substitutes. \textless{}/think\textgreater{} \textbackslash\{\}boxed\{C\}} $\to$ \texttt{C}
\end{quote}

\noindent\textbf{Overconfident-wrong $\to$ calibrated-right} (\texttt{mmlu-12923}; gold C
, M4 state overconfident wrong $\to$ nonconfident right). The baseline burns its budget re-asserting; the steered trace weighs all four options tersely and lands on the same letter with calibrated confidence.
\begin{quote}\scriptsize\emph{Question:} The best conclusion that can be drawn about the four-day ( compressed) workweek is that it has: Options: A. long-lasting positive effects on attitudes and productivity. B. long-lasting negative effects on attitudes and p\\[2pt]
\textbf{before:} \texttt{\textless{}think\textgreater{}The question asks for the best conclusion about the four-day workweek. We need to consider the potential effects on attitudes and productivity. There's no mention of specific evidence or studies, so we need to consider the general impact of a compressed workweek. \textless{}/think\textgreater{} \textless{}think\textgreater{}Option A is too strong. We don't have enough information to say it has long-lasting positive effects on attitudes and productivity. \textless{}/think\textgreater{} \textless{}think\textgreater{}Option B is too strong. We don't have enough information to say it has long-las} $\to$ \texttt{D}\\[2pt]
\textbf{after:} \texttt{\textless{}think\textgreater{}The question asks for the best conclusion about the four-day workweek. We need to consider the potential effects on attitudes and productivity. Option A is too strong, as we don't have enough information to guarantee long-lasting positive effects. Option B is too strong, as we don't have enough information to guarantee long-lasting negative effects. Option C is a possibility, but we don't know if it's the best conclusion. Option D is a possibility, but we don't know if it's the best conclusion. \textless{}/think\textgreater{} \textbackslash\{\}} $\to$ \texttt{D}
\end{quote}

\noindent\textbf{Confident-right, gated but spared} (\texttt{mmlu-5125}; gold B
, M4 state confident right $\to$ confident right). A false-flagged correct answer: the steered reasoning changes wording, the answer improves/stays correct, and the confident-right state survives.
\begin{quote}\scriptsize\emph{Question:} Which one of the following best states Ivan Pavlov's accidental discovery? Options: A. While experimenting with rats, Pavlov discovered that if he repeatedly placed pellets of food in one side of a "T-maze," the rats wou\\[2pt]
\textbf{before:} \texttt{\textless{}think\textgreater{}The question asks for Pavlov's accidental discovery. Let's look at the options and see which one describes an accidental discovery. Option B describes a learned response to a stimulus, which is not an accidental discovery. Option C describes a cat learning to press a lever, which is a learned behavior. Option D describes a pigeon learning to associate a key with food. Option A describes a rat learning to run to a specific side of a maze. This is a classic example of classical conditioning. Therefore,} $\to$ \texttt{A}\\[2pt]
\textbf{after:} \texttt{\textless{}think\textgreater{}The question asks for the accidental discovery of Pavlov. Option B is the best choice because it describes Pavlov's accidental discovery of the dog's salivation response to the sight and sound of the attendant who brought them food. \textless{}/think\textgreater{} \textbackslash\{\}boxed\{B\}} $\to$ \texttt{B}
\end{quote}
